\documentclass{article}
\usepackage{hyperref}
\usepackage{Style}

\nocite{*} % Comentar si quiero citar
%\addbibresource{bibliografia.bib} % Quitar el comentado si quiero usar bibliografia

\begin{document}

\begin{minipage}{2.5cm}
    \includegraphics[width=2cm]{imagen_puc.jpg}
\end{minipage}
\begin{minipage}{12cm}
    {\sc Pontificia Universidad Católica de Chile\\
    Facultad de Matemáticas\\
    Departamento de Matemática\\
    Profesor: Manuel Cabezas -- Ayudante: Benjamín Mateluna}
\end{minipage}
\vspace{2ex}

{\centerline{\bf Introducción a la Geometría - MAT1304}
\centerline{\bf Ayudantía 8}}
\centerline{\bf 7 de abril de 2026}

\vspace{1cm}
\noindent\textbf{Problema 1.} Sea $\triangle ABC$ un triángulo cualquiera y puntos $D,E,F$ en los
lados $\overline{BC}, \overline{AB}$ y $\overline{AC}$ respectivamente. Muestre que las 
circunferencias circunscritas a los triángulos $\triangle AEF,\triangle BED$ y $\triangle CFD$ son
concurrentes.

\vspace{5mm}
\noindent\textbf{Problema 2.} Sean $A$ y $B$ puntos en la circunferencia. Sean $P$ y $Q$ puntos en
el arco $BA$ y sea $X$ el punto en $AB$ donde la bisectriz del ángulo 
$\angle APB$ corta a la circunferencia. Demuestre que $\overline{QX}$ es la bisectriz del ángulo 
$\angle AQB$.

\vspace{5mm}
\noindent\textbf{Problema 3.} Sean $\overline{AC}$ y $\overline{BD}$ cuerdas perpendiculares en 
una circunferencia que se intersectan en el punto $G$. En $\triangle ADG$ la altura en $G$ 
intersecta a $\overline{AD}$ en $E$ y si la extendemos intersecta a $\overline{BC}$ en $P$. 
Muestre que $P$ es punto medio de $\overline{BC}$.

\vspace{5mm}
\noindent\textbf{Problema 4.} Construir $\triangle ABC$ dados lado $a$, lado $b$ y altura $h_{B}$.

\vspace{1cm}
\noindent\textbf{\underline{Ejercicios Propuestos:}}

\vspace{5mm}
\noindent\textbf{Extra 1.} Mostrar que un polígono de 5 lados convexo con vértices consecutivos
$A,B,C,D$ y $E$ está inscrito en una circunferencia si y solo si
\begin{equation*}
    \angle EBC+\angle EDC=180^{\circ} \hhtext{y} \angle DAB+\angle DCB= 180^{\circ}
\end{equation*}

\vspace{5mm}
\noindent\textbf{Extra 2.} Dada una recta $L$ considere los puntos $M$ y $N$ al mismo lado de
$L$. Construya el punto $C$ en $L$ tal que el ángulo que forma $L$ con $\overleftrightarrow{MC}$
es congruente con el ángulo determinado por $L$ con $\overleftrightarrow{NC}$.

%\printbibliography % Quitar el comentado si quiero usar bibliografia

\end{document}